% !PRINTED 653
\weberpage{653}

Diese beiden Zahlen gestatten, wie aus (7) unmittelbar hervorgeht, die Substitutionen $(x_0, x_\varkappa)$, und sind also Zahlen des Körpers $\Omega_m$.

Ist $n$ durch $q$ nicht theilbar, so kann wegen der Irreducibilität der Kreistheilungsgleichung keine von den Resolventen $(r^n, x_0)$ verschwinden, wenn sie nicht alle verschwinden; denn die Gleichung
\[
(r, x_0)^m = 0
\]
wäre eine Gleichung für $r$ mit rationalen Co\"efficienten, die alle Substitutionen $(r, r^n)$ gestattet.

Ebenso kann keine der Resolventen $(r^s, x_0)$ verschwinden, wenn nicht die ganze Schaar verschwindet, in der $s$ einen bestimmten grössten gemeinschaftlichen Theiler mit $m$ hat.

Wenn alle $(r, x_0)$ verschwinden, so ist die Formel (9) nicht anwendbar. Dann können wir aber die Resolventen $(r^q, x_0)$ direct betrachten, die zu einem Körper vom Grade $m q^{-1}$ gehören, der aus
\[
y_0 = x_0 + x_q + x_{2q} + \cdots + x_{m-q}
\]
hervorgeht, und wenn auch diese Grössen verschwinden, so gehen wir zu den Resolventen $(r^{q^2}, x_0)$ über u.\,s.\,f.

Wenn aber $(r, x_0)$ nicht verschwindet, so lassen sich durch die Formel (9) alle $(r^s, x_0)$ rational durch $(r, x_0)$ und durch Kreistheilungszahlen darstellen.

Demnach ist es für unseren Zweck ausreichend, wenn wir beweisen können, dass die \gesperrt{nicht verschwindenden Resolventen} $(r, x_0)$ \gesperrt{Kreistheilungszahlen sind}.

Lassen wir $n$ ein volles System incongruenter, durch $q$ nicht theilbarer Zahlen durchlaufen, so ist
\[
\sum n \equiv 0 \ (\mathrm{mod}\ m),
\]
denn die Zahlen $n$ zerfallen in Paare einander zu $m$ ergänzender Zahlen. Daraus folgt nach (7), dass das Product
\weq{10}{\prod^n (r^n, x_0) = a}
eine Zahl in $\Omega_m$ ist. Da aber diese Zahl sich überdies nicht ändert, wenn irgend eine der Substitutionen $(r, r^n)$ darin ausgeführt wird, so ist $a$ eine \gesperrt{rationale Zahl}.
